\emph{(i).} By \claimref{tower.ladder-pairs-satisfy-car-vectors} the ladder pairs satisfy $v_{j}w_{k} + w_{k}v_{j} = \delta_{jk}$, $v_{j}v_{k} + v_{k}v_{j} = 0 = w_{j}w_{k} + w_{k}w_{j}$. The adjoint is the linear extension of the inverse on blades; it is an anti-automorphism with $e_{j}^{\dagger} = e_{j}$, $f_{j}^{\dagger} = -f_{j}$ (the group of signed blades of \claimref{synthesis.signed-directions-and-the-integer-clifford-algebra}, on which the inverse reverses products), so $v_{j}^{\dagger} = w_{j}$. Each $p_{j} = w_{j}v_{j}$ is a self-adjoint idempotent, the $p_{j}$ commute, $P_{0} = \prod_{j}p_{j}$ is a self-adjoint idempotent, $P_{0} \ne 0$ (since $p_{j} = \tfrac12(1 + e_{j}f_{j})$ and the scalar part of $P_{0}$ is $2^{-\ell}$), and $w_{j}P_{0} = 0$ because $w_{j}w_{j}v_{j} = 0$. Moving the $w$'s of $s_{I}^{\dagger} = P_{0}w_{i_{m}}\cdots w_{i_{1}}$ to the right through $s_{J}$ with the relations gives $s_{I}^{\dagger}s_{J} = \delta_{IJ}P_{0}$. The left ideal $\Cl_{\Q}(\ell,\ell)P_{0}$ is spanned by the $2^{\ell}$ elements $s_{I}$ (a word in the $v$'s and $w$'s applied to $P_{0}$ is, after ordering, zero or $\pm$ some $s_{I}$), so $xs_{J} = \sum_{I}m(x)_{IJ}s_{I}$ defines a homomorphism $m$ to $\Mat_{2^{\ell}}(\Q)$, and $s_{I}^{\dagger}xs_{J} = m(x)_{IJ}P_{0}$. Then $m(x^{\dagger})_{IJ}P_{0} = s_{I}^{\dagger}x^{\dagger}s_{J} = (s_{J}^{\dagger}xs_{I})^{\dagger} = m(x)_{JI}P_{0}$. The images $m(v_{j})$, $m(w_{j})$ are the creation and annihilation matrices of the $2^{\ell}$ basis vectors, which generate the full matrix algebra (the matrix unit carrying $s_{J}$ to $s_{I}$ is the image of $s_{I}s_{J}^{\dagger}$, since $s_{I}s_{J}^{\dagger}s_{K} = \delta_{JK}s_{I}$); both algebras have dimension $4^{\ell}$, so $m$ is an isomorphism.

\emph{(ii).} By clause (ii) of \claimref{synthesis.positive-rationals-are-the-squares-of-the-even-subalgebra-of-cl-three} the positive rational $r \cdot r$ is a sum of four rational squares; put them in the four places of $T$ to get $r' \in \Q^{T}$ with $r' \cdot r' = r \cdot r$. Since $S$ has an index outside $T$, $r \notin \Q^{T}$ and $u = r - r' \ne 0$; the form is positive definite, so $u \cdot u \ne 0$ and $g_{u}(x) = -uxu^{-1} = x - 2\tfrac{x \cdot u}{u \cdot u}u$ is the reflection of \claimref{std-algebra.reflection}, an isometry with rational coefficients. From $r \cdot r = r' \cdot r'$, $u \cdot u = 2(r' \cdot r' - r \cdot r') = -2\,r' \cdot u$, so $g_{u}(r') = r' + u = r$. For $x \in \Q^{T}$ write $c(x) = 2(x \cdot u)/(u \cdot u)$, a linear functional on $\Q^{T}$; then $g_{u}(x) = x - c(x)(r - r')$, whose components are
\[
g_{u}(x)_{t} = x_{t} - c(x)(r_{t} - r'_{t}) =: \lambda_{t}(x) \ \ (t \in T), \qquad g_{u}(x)_{i} = -c(x)\,r_{i} \ \ (i \notin T).
\]
So every vector of $W = g_{u}(\Q^{T})$ has support in $T \cup S = S$. The five functionals $\lambda_{t}$ ($t \in T$) and $c$ on $\Q^{T}$ are nonzero: $c(r') = -1$ and $\lambda_{t}(r') = r_{t} \ne 0$, because $g_{u}(r') = r$ and $t \in S$. Their kernels are therefore hyperplanes of $\Q^{T}$, and $g_{u}$ carries them to hyperplanes of $W$: the kernel of $\lambda_{t}$ to the vectors of $W$ with vanishing $t$-th component, which have support in $S \setminus \{t\}$, and the kernel of $c$ to the vectors $g_{u}(x) = x$, which lie in $\Q^{T}$. If all five functionals vanish at $x$ then $c(x) = 0$ and $x_{t} = \lambda_{t}(x) = 0$ for all $t$, so $x = 0$: the five functionals span the dual of the four-dimensional space $\Q^{T}$, so four of them are linearly independent, and independent functionals have distinct kernels.

\emph{A lemma.} Let $q$ be a quadratic form on a rational vector space $U$ of finite dimension, with values in $V$ (a function $q(x) = \beta(x,x)$ with $\beta$ symmetric and bilinear), vanishing on three distinct hyperplanes $\ker\alpha_{1}$, $\ker\alpha_{2}$, $\ker\alpha_{3}$. Then $q = 0$. Proof: choose coordinates with $\alpha_{1}$ the first coordinate; $q$ has no monomial free of $x_{1}$, so $q = \alpha_{1}\cdot\gamma$ with $\gamma$ linear with values in $V$. On $\ker\alpha_{2}$, which is not $\ker\alpha_{1}$, the product $\alpha_{1}\gamma$ vanishes while $\alpha_{1}$ is a nonzero linear functional there, so $\gamma$ vanishes on $\ker\alpha_{2}$ (a linear map vanishing off a proper subspace vanishes), that is $\gamma = \alpha_{2}v$ for some $v \in V$. So $q = \alpha_{1}\alpha_{2}v$, and a point of $\ker\alpha_{3}$ outside $\ker\alpha_{1} \cup \ker\alpha_{2}$, which exists because a vector space over $\Q$ is not the union of two proper subspaces, gives $v = 0$.

\emph{(iii).} Let $\mu$ be an additive weight and put $q(r) = (r \cdot r)\,\mu(\Pi_{r})$ for $r \ne 0$, $q(0) = 0$. Since $\Pi_{\lambda r} = \Pi_{r}$, $q(\lambda r) = \lambda^{2}q(r)$.

\emph{Step 1: on the image of a coordinate four-space under an isometry, $q$ is a quadratic form.} Let $T$ be four indices, $\iota : \Q^{4} \to \Q^{T}$ the coordinate identification and $g$ an orthogonal matrix with rational entries (a product of reflections and coordinate permutations will do). For a projector $E$ of $\Mat_{4}(\Q)$ let $\hat{E} = g\,\iota E\iota^{T}g^{T}$, a projector of $\Mat_{n}(\Q)$, with $\hat{E}\hat{E}' = g\,\iota EE'\iota^{T}g^{T}$. So $E \mapsto \mu(\hat{E})$ is an additive weight on the projectors of $\Mat_{4}(\Q)$ with values in $V$, and by clause (iii) of \claimref{synthesis.additive-weights-on-the-rational-projectors-of-two-modes-are-linear} there are $\rho'_{ab} = \rho'_{ba} \in V$ with $\mu(\widehat{\Pi_{x}}) = \sum\rho'_{ab}x_{a}x_{b}/(x \cdot x)$. Now $\widehat{\Pi_{x}} = \Pi_{g\iota x}$ and $g\iota x \cdot g\iota x = x \cdot x$, so $q(g\iota x) = \sum\rho'_{ab}x_{a}x_{b}$: the function $x \mapsto q(g\iota x)$ is a quadratic form on $\Q^{4}$.

\emph{Step 2: the form $\rho$.} Put $\rho_{ii} = q(e_{i})$ and $\rho_{ij} = \tfrac12\bigl(q(e_{i} + e_{j}) - q(e_{i}) - q(e_{j})\bigr)$ for $i \ne j$, and $D(r) = q(r) - \sum\rho_{ij}r_{i}r_{j}$. For four indices $T$ (and $n \ge 4$ puts every pair of indices in such a $T$), Step 1 with $g = 1$ says that $q$ restricted to $\Q^{T}$ is a quadratic form; a quadratic form is determined by its values on the $e_{i}$ and the $e_{i} + e_{j}$, so on $\Q^{T}$ it equals $\sum_{i,j \in T}\rho_{ij}r_{i}r_{j}$. Hence $D(r) = 0$ for every $r$ whose support has at most four indices. By Step 1, $D \circ g \circ \iota$ is a quadratic form on $\Q^{4}$ for every $T$ and every rational orthogonal $g$, being the difference of two.

\emph{Step 3: induction on the support.} Suppose $D$ vanishes on all vectors whose support has at most $k$ indices, $k \ge 4$, and let $r$ have support $S$ with $k + 1$ indices. Take $T$, $r'$, $u$ and $W = g_{u}(\Q^{T})$ as in (ii). The quadratic form $D \circ g_{u}$ on $\Q^{T}$ vanishes on the kernel of each $\lambda_{t}$, whose image consists of vectors with support in $S \setminus \{t\}$, at most $k$ indices, and on the kernel of $c$, whose image lies in $\Q^{T}$, four indices. By (ii) at least four of these hyperplanes are distinct, so by the lemma $D \circ g_{u} = 0$ on $\Q^{T}$, and $D(r) = D(g_{u}(r')) = 0$. So $D = 0$: $\mu(\Pi_{r}) = \sum\rho_{ij}r_{i}r_{j}/(r \cdot r)$ for all $r$. A projector $E$ is the orthogonal projection onto a rational subspace, which has an orthogonal basis $r_{1}, \dots, r_{m}$ of rational vectors (Gram-Schmidt needs no square roots when the vectors are not normalised), $E = \sum_{a}\Pi_{r_{a}}$ with the $\Pi_{r_{a}}$ pairwise orthogonal, and by additivity $\mu(E) = \sum_{a}\sum_{i,j}\rho_{ij}(\Pi_{r_{a}})_{ij} = \sum_{i,j}\rho_{ij}E_{ij}$. Uniqueness: $\rho_{ij}$ is recovered from the values of $\mu$ at $\Pi_{e_{i}}$ and $\Pi_{e_{i} + e_{j}}$.

\emph{(iv).} Let $\mu$ be as stated and $\rho$ the symmetric rational matrix of (iii), so $\rho(r,r) := \sum\rho_{ij}r_{i}r_{j} = (r \cdot r)\mu(\Pi_{r}) \ge 0$ for every rational $r$ and $\sum_{i}\rho_{ii} = \mu(1) = 1$. Put $\omega(x) = \sum_{i,j}\rho_{ij}\,m(x)_{ij}$ with $m$ of (i). Then $\omega(1) = 1$, $\omega(x^{\dagger}) = \omega(x)$ because $m(x^{\dagger}) = m(x)^{T}$ and $\rho$ is symmetric, and $\omega(xx^{\dagger}) = \sum_{k}\rho(y_{k}, y_{k}) \ge 0$ with $y_{k}$ the columns of $m(x)$. It restricts to $\mu$ on the projectors by (iii). A state restricts to an additive weight with non-negative values ($E = EE^{\dagger}$) and total $1$, and two states with the same restriction have the same symmetric matrix by the uniqueness in (iii), hence agree on symmetric elements and, by $\omega(x^{\dagger}) = \omega(x)$, on all elements.

\emph{(v).} Let $u = \varphi(x)$, $w = \varphi(y)$. Then $(u \cdot u)(w \cdot w) - (u \cdot w)^{2} = c^{2}\bigl((x \cdot x)(y \cdot y) - (x \cdot y)^{2}\bigr)$. Read $x, y \in \Q^{4}$ as even elements of $\Cl_{\Q}(3,0)$ (\claimref{std-algebra.quaternions}), with $N$ the square of clause (i) of \claimref{synthesis.positive-rationals-are-the-squares-of-the-even-subalgebra-of-cl-three}: $N(x) = x \cdot x$, the scalar part of $x\tilde{y}$ is $x \cdot y$, and $N(x\tilde{y}) = N(x)N(y)$. So $(x \cdot x)(y \cdot y) - (x \cdot y)^{2}$ is $N(x\tilde{y})$ minus the square of its scalar part, the sum of the squares of the three bivector components of $x\tilde{y}$; times $c^{2}$ it is still a sum of three rational squares. That the left side is $-(u \wedge w)^{2}$ is the expansion $(uw)(wu) = (u \cdot w + u \wedge w)(u \cdot w - u \wedge w)$ in $\Cl_{\Q}(n,0)$, whose left side is $(u \cdot u)(w \cdot w)$. For $u = e_{1}$, $w = e_{2} + e_{3} + e_{4} + 2e_{5}$: $u \cdot u = 1$, $w \cdot w = 7$, $u \cdot w = 0$, so the value is $7$, which is not a sum of three rational squares by clause (iv) of the cited result. The same pair lies in $\Q^{n}$ for every $n \ge 5$.
